%
% File: methodology
%
\chapter{Methodology}
\label{chap:method}
\setlength{\parskip}{1em}
%
\section{Research Design}

This study employs a systematic literature review methodology to synthesise current knowledge on AI-generated visual explanations as a pedagogical tool for learners with low language proficiency or learning disabilities. A systematic review is appropriate because the field is fragmented across multiple sub-disciplines --- educational technology, special education, cognitive psychology, and human--computer interaction --- and because the rapid pace of development in generative AI requires periodic consolidation of evidence. The review follows the principles of the Preferred Reporting Items for Systematic Reviews and Meta-Analyses (PRISMA) statement, with adaptations appropriate to a primarily qualitative thematic synthesis. Specifically, the review uses PRISMA's four-phase approach (identification, screening, eligibility, inclusion) for source selection, but uses thematic analysis rather than meta-analytic statistical pooling for synthesis, given the heterogeneity of the included studies in design, population, and outcome measures.

\section{Search Strategy}

The search was conducted between September and December 2025 across five academic databases selected for their coverage of education, psychology, and information science: Scopus, Web of Science, ERIC, IEEE Xplore, and Google Scholar. The search strategy combined three concept blocks using Boolean operators. Block A targeted the technology: ``generative AI'' OR ``large language model'' OR ``multimodal AI'' OR ``AI-generated image'' OR ``visual generation''. Block B targeted the pedagogical context: ``education'' OR ``learning'' OR ``instruction'' OR ``pedagogy'' OR ``multimedia learning''. Block C targeted the population: ``learning disabilit*'' OR ``language proficiency'' OR ``English language learner'' OR ``special education'' OR ``dyslexia'' OR ``autism'' OR ``ADHD''. The three blocks were combined as A AND B AND C, with database-specific syntax adjustments. The search was limited to peer-reviewed journal articles, edited book chapters from academic publishers, and conference proceedings from established venues. Date limits were 2023--2025, reflecting the post-GPT-4 era of multimodal generative AI. The search was conducted in English, which constitutes a limitation acknowledged in Chapter~\ref{chap:discussion}.

\section{Inclusion and Exclusion Criteria}

Studies were included if they met all of the following criteria. First, the study addressed AI tools capable of generating, manipulating, or contextualising visual content, defined broadly to include diagrams, illustrations, animations, augmented reality overlays, and image-based augmentative communication. Second, the study addressed an educational context, including formal schooling, higher education, and structured non-formal learning. Third, the study either explicitly addressed learners with low language proficiency or learning disabilities, or its findings were directly transferable to those populations as evidenced by the use of relevant theoretical frameworks (CTML, CLT, universal design for learning). Fourth, the study reported either empirical findings, a systematic synthesis of empirical findings, or a theoretically grounded framework with explicit pedagogical implications.

Studies were excluded if they addressed AI in education without any visual or multimodal component (for example, text-only chatbot studies); addressed visual technology in education without an AI component (for example, conventional video learning); were limited to technical descriptions of AI systems without pedagogical analysis; were published in non-peer-reviewed venues; or were not available in English. Opinion pieces, editorials, and short commentaries were also excluded.

\section{Selection Process}

The initial database search returned approximately 340 records. After duplicate removal across databases, 287 unique records remained. Title and abstract screening reduced this to 42 records assessed for full-text eligibility. Full-text review excluded 32 records, primarily because they addressed AI in education without a visual generation component ($n=18$), addressed visual learning without AI ($n=7$), or were not directly relevant to the target populations ($n=7$). The final corpus comprised ten studies, listed in the references and summarised in Table~\ref{tab:studies} (Chapter~\ref{chap:findings}). The selection process is summarised in the PRISMA flow diagram presented as Figure~\ref{fig:prisma}.

\begin{figure}[H]
\centering
\begin{tabular}{|c|}
\hline
\textbf{Identification:} 340 records identified through database searching \\
\hline
$\downarrow$ \\
\hline
\textbf{Screening:} 287 unique records after duplicate removal \\
\hline
$\downarrow$ \\
\hline
\textbf{Eligibility:} 42 records assessed for full-text eligibility \\
\hline
$\downarrow$ \\
\hline
\textbf{Inclusion:} 10 studies included in the final synthesis \\
\hline
\end{tabular}
\caption{PRISMA-style flow diagram of the source selection process.}
\label{fig:prisma}
\end{figure}

\section{Data Extraction}

A structured data extraction form was developed to ensure consistent treatment of each included study. The form captured: bibliographic information (authors, year, journal); study type (systematic review, empirical study, theoretical framework, qualitative study); target population (specific disability or language-proficiency profile, age group, educational level); AI tool or application examined; visual modality involved (static image, diagram, animation, AR, multimodal); theoretical framework cited; methodology; key findings; and reported limitations. Extraction was conducted by the authors and entered into a single spreadsheet, with each row representing one study. For studies that were themselves systematic reviews, the extraction recorded summary-level findings rather than re-extracting from primary sources, in order to avoid double-counting.

\section{Synthesis Approach}

Synthesis followed Braun and Clarke's six-phase thematic analysis procedure, adapted for documentary rather than interview data. After familiarisation with each study, initial codes were generated inductively from the extracted data, with attention both to manifest content (what the studies reported) and latent content (theoretical assumptions, recurring concerns). Codes were then collated into candidate themes, which were reviewed against the full dataset to ensure internal coherence and external distinctiveness. Four themes were finalised: (i)~typology of AI-generated visual explanations; (ii)~alignment with multimedia learning theory; (iii)~evidence of pedagogical efficacy; and (iv)~implementation barriers and ethical concerns. Each theme is presented in Chapter~\ref{chap:findings} with supporting evidence from the included studies.

\section{Quality and Limitations}

Quality appraisal followed the criteria appropriate to each study type: systematic reviews were appraised against AMSTAR-2 criteria, empirical studies against the Mixed Methods Appraisal Tool, and theoretical papers against criteria for conceptual rigour and engagement with prior literature. All ten included studies met the threshold for inclusion. The principal limitations of the review --- language restriction, narrow date range, single-team extraction --- are acknowledged in Chapter~\ref{chap:discussion}. The methodological choices made here represent a balance between systematic rigour and feasibility within the scope of a course-bounded review.
